\section{A mixed determinant construction and a staircase gate}

\subsection*{Finding F-017: split early variables and later degree}
Let \(m\geq2\), \(n>m\), and \(\lambda\vdash nd\), with
\(r=\ell(\lambda)\). Fix \(2\leq k\leq r\) and an integer
\(s\geq1\) such that \(\lambda_k\leq s\). Write
\(R=r-k+1\). Choose nonnegative integers
\[
                   p+q=k-2,\qquad u+v=s-1,
\]
and set
\begin{equation}\label{eq:F017size}
 \mathcal S_m(p,q,u,v;R)=
 \binom{m+p}{p}\binom{R+u}{u}
 +\binom{m+q}{q}\binom{R+v}{v}-1.
\end{equation}
Then
\begin{equation}\label{eq:F017compare}
 n\geq\mathcal S_m(p,q,u,v;R)
       \quad\Longrightarrow\quad
                  q^m_\lambda(d[n])\leq K_n(\lambda).
\end{equation}
The graph establishing the bound mixes a split of the early
variables with a split of the degree in the later variables.
Its coefficient edge carries one later variable when there
is one. This saves a full degree from the two stored parts.

\paragraph{The universal determinant identity.}
Let \(L\), \(z_1,\ldots,z_{k-2}\), and
\(y_1,\ldots,y_R\) be variables. Suppose that a homogeneous
degree-\(m\) form has degree at most \(s\) in the \(y\)'s:
\[
 f=\sum_{|\alpha|+|\gamma|\leq m,\ |\gamma|\leq s}
              c_{\alpha,\gamma}z^\alpha y^\gamma
                                      L^{m-|\alpha|-|\gamma|}.
\]
There is a homogeneous linear matrix with
\(S=\mathcal S_m(p,q,u,v;R)\) vertices and determinant
\begin{equation}\label{eq:F017identity}
 \det M=
 \sum_{|\alpha|+|\gamma|\leq m,\ |\gamma|\leq s}
      c_{\alpha,\gamma}z^\alpha y^\gamma
                                      L^{S-|\alpha|-|\gamma|}.
\end{equation}
Every exponent actually occurring on the right is nonnegative.
Thus, for \(n\geq\max\{m,S\}\), adjoining
\(L I_{n-S}\) represents \(L^{n-m}f\) as an
\(n\)-by-\(n\) homogeneous linear determinant.
The coefficients of this identity are integral polynomials
in the parameters \(c_{\alpha,\gamma}\).

\begin{proof}
Split the early variables \(z\) into disjoint sets of sizes
\(p\) and \(q\). The first tree has as vertices pairs
\((\alpha_A,\gamma_A)\), where \(\alpha_A\) is a
multi-index in the first set with \(|\alpha_A|\leq m\)
and \(\gamma_A\) is a multi-index in \emph{all} \(R\)
later variables with \(|\gamma_A|\leq u\).
The second tree has the analogous pairs
\((\alpha_B,\gamma_B)\), with the other early variables
and \(|\gamma_B|\leq v\).
Use a fixed parent rule which decreases the last nonzero
coordinate of a pair. The sets are coordinatewise downward
closed, so every root path remains present and is unique.
Orient the first tree away from its root and the second
toward its root. Merge the roots and no other vertices.
In particular, the two copies of a later-variable multi-index
are distinct vertices, despite having the same variable labels.
Their total number is exactly \(S\) in \eqref{eq:F017size}.

All tree edges have labels equal to the negative of the
variable they add or remove. For a monomial with \(\gamma=0\),
insert an edge from \((\alpha_A,0)\) to
\((\alpha_B,0)\), labeled \(c_{\alpha,0}L\).
For \(\gamma\ne0\), choose one \(y_j\) with
\(\gamma_j>0\) by a fixed rule. Split
\[
              \gamma-e_j=\gamma_A+\gamma_B,
                     \qquad |\gamma_A|\leq u,
                            \quad |\gamma_B|\leq v.
\]
This is possible because \(|\gamma|-1\leq s-1=u+v\);
taking a fixed initial segment of its sorted variable word
gives a unique choice. Add the edge from
\((\alpha_A,\gamma_A)\) to
\((\alpha_B,\gamma_B)\), labeled
\(c_{\alpha,\gamma}y_j\).
Parallel labels are added in the relevant matrix entry.
Put an \(L\)-loop at every noncentral vertex. The central
diagonal consists of its coefficient-loop labels only.

Every non-loop cycle passes through the central vertex and
uses exactly one coefficient edge. Hence a cycle cover chooses
one coefficient-edge summand and its two unique tree paths,
with loops at all unused vertices.
For \(\gamma=0\), that cycle has \(|\alpha|+1\) edges,
of which \(|\alpha|\) are negative tree edges. The
permutation sign cancels those signs and gives
\(c_{\alpha,0}z^\alpha L^{S-|\alpha|}\).
For \(\gamma\ne0\), the paths have
\(|\alpha|+|\gamma|-1\) tree edges and the cycle has
\(|\alpha|+|\gamma|\) edges. Its sign again cancels
the tree signs; the coefficient edge supplies the removed
\(y_j\). The contribution is
\(c_{\alpha,\gamma}z^\alpha y^\gamma
L^{S-|\alpha|-|\gamma|}\).
This includes a central coefficient loop when the paths
are empty. Since each cycle uses distinct vertices, all
displayed exponents are nonnegative. Summing proves
\eqref{eq:F017identity} and the claimed padding identity.
Writing the result term by term avoids an unjustified
negative power when \(S<m\).
\end{proof}

\begin{proof}[Proof of the actual multiplicity comparison]
If \(q^m_\lambda=0\), the assertion is immediate. Otherwise
\(r\leq\min\{d,m^2\}\), so the natural highest-weight
space is supported on an \(r\)-dimensional coordinate subspace.
Let \(w\) be any highest-weight vector of weight \(\lambda\)
in the determinant ideal. A restricted generic translate
of the internally padded permanent is \(L^{n-m}f\),
where \(f\) has degree \(m\). On the nonempty open set
\(L(e_1)\ne0\), the upper-unipotent change of variables
from M-009 aligns \(L\) with \(x_1\) without changing
the vanishing of \(w\).

The support theorem F-016 says that no coefficient generator
used by \(w\) has total exponent greater than \(\lambda_k\)
in \(x_k,\ldots,x_r\). Delete all monomials of \(f\)
with larger degree in those variables and call the result
\(f_{\leq\lambda_k}\). Padding changes only the exponent
of \(x_1\), so
\[
       w(x_1^{n-m}f)=w(x_1^{n-m}f_{\leq\lambda_k}).
\]
The retained form has later-variable degree at most
\(\lambda_k\leq s\). Apply \eqref{eq:F017identity}
with \(z=(x_2,\ldots,x_{k-1})\) and
\(y=(x_k,\ldots,x_r)\). Under the size hypothesis, its
padded form is a determinant and \(w\) vanishes on it.
Thus every such \(w\) vanishes on a dense open subset
of the permanent orbit and hence on its reduced closure.
The highest-weight ideal multiplicity spaces satisfy
\(J_{\Omega_n,\lambda}\subseteq J_{Z_{n,m},\lambda}\).
This is an inclusion for every copy; taking quotient dimensions
proves \eqref{eq:F017compare}.
\end{proof}

\paragraph{A staircase in the BIP numerical range.}
Suppose \(n\geq m^{25}\). Every actual multiplicity
obstruction satisfies the twenty-five inequalities
\begin{equation}\label{eq:F017BIP}
                \lambda_{54-2s}\geq s+1
                               \qquad(1\leq s\leq25).
\end{equation}
In particular,
\[
       \lambda_4\geq26,\quad\lambda_6\geq25,\quad
       \lambda_8\geq24,\quad\ldots,\quad\lambda_{52}\geq2.
\]
Odd-indexed consequences follow from monotonicity of the rows.
For example \(\lambda_3\geq26\) and \(\lambda_{27}\geq14\).

\begin{proof}
By M-008 an actual obstruction already has \(r\geq52\),
so every index \(k=54-2s\) in \eqref{eq:F017BIP} exists.
Suppose \(\lambda_k\leq s\). Choose
\(u=\lfloor(s-1)/2\rfloor\),
\(v=\lceil(s-1)/2\rceil\), and
\[
                     p=25-2u,\qquad q=25-2v.
\]
Then \(0\leq u,v\leq12\), \(p,q\geq1\), and
\(p+q=k-2\). Each term of \eqref{eq:F017size}
has total power budget \(p+2u=q+2v=25\).

Here are uniform estimates including all small parameters.
For \(a\geq0\), \(m\geq2\), the product formula gives
\[
 \frac{\binom{m+a}{a}}{m^a}
       \leq A_a:=\frac{\binom{a+2}{2}}{2^a}.
\]
Since \(R\leq r-1\) and \(r\geq2\), similarly
\[
 \frac{\binom{R+j}{j}}{r^j}
       \leq\frac{\binom{r-1+j}{j}}{r^j}
       \leq B_j:=\frac{j+1}{2^j}.
\]
Both inequalities follow by observing that each factor
\(1+i/m\), or \(1+i/r\), decreases as its denominator
increases. The formulas include \(a=0\) and \(j=0\).
The sequence \(A_a\) is at most \(3/2\) and decreases
for \(a\geq1\); \(B_j\leq1\) and decreases for \(j\geq1\).
For \(0\leq j\leq12\), with \(a=25-2j\), this gives
\[
 A_a B_j\leq
 \begin{cases}
 55/512,&j\leq8,\quad a\geq9,\\
 15/512,&j\geq9,
 \end{cases}
\]
because \(A_9=55/512\) and \(B_9=10/512\).
Using \(r\leq m^2\), each term in the size formula is
therefore at most \((55/512)m^{25}\), and
\[
                 \mathcal S_m
                   \leq\frac{110}{512}m^{25}-1<m^{25}\leq n.
\]
The comparison theorem forbids the alleged obstruction,
which proves \eqref{eq:F017BIP}.
\end{proof}

\paragraph{A staircase in the IP numerical range.}
If \(n>3m^4\), every actual multiplicity obstruction satisfies
\begin{equation}\label{eq:F017IP}
              \lambda_{12-2s}\geq s+1
                                     \qquad(1\leq s\leq5),
\end{equation}
or explicitly
\[
       \lambda_2\geq6,\quad\lambda_4\geq5,\quad
       \lambda_6\geq4,\quad\lambda_8\geq3,\quad\lambda_{10}\geq2.
\]

\begin{proof}
M-008 gives \(r\geq10\) for an actual obstruction.
If \(k=12-2s\) and \(\lambda_k\leq s\), choose
\(u=\lfloor(s-1)/2\rfloor\),
\(v=\lceil(s-1)/2\rceil\),
\(p=4-2u\), and \(q=4-2v\).
These are nonnegative and satisfy all required sums,
with \(p+2u=q+2v=4\).
The estimates just proved give \(A_p,A_q\leq3/2\)
and \(B_u,B_v\leq1\). Thus
\[
                   \mathcal S_m\leq3m^4-1<n.
\]
Again \eqref{eq:F017compare} excludes the obstruction.
\end{proof}

\paragraph{The corresponding full ambient equality.}
Assume \(r\leq\min\{d,n^2\}\). For \(b=\lambda_2\), replace \(m\) by \(b\) in
\eqref{eq:F017size}, take \(s\geq\lambda_k\), and assume
\(n\geq\max\{b,\mathcal S_b\}\).
Then the same coefficient-space construction, F-014, and
the restriction-and-padding proof of F-015 give
\[
                         K_n(\lambda)=a_\lambda(d[n]).
\]
Indeed all stabilized highest-weight vectors use only the
degree-\(b\) coefficient space with later-variable degree
at most \(\lambda_k\). The differential padding composition
is diagonal and invertible on this monomial space, so every
ambient copy remains independent. This full-equality variant
has a different size parameter from the direct degree-\(m\)
comparison above. The staircase conclusions use the latter
and make no assumption that \(n\geq\lambda_2\).

\paragraph{Dictionary and remaining gate.}
The coefficient support is detected by exterior matrix
coefficients of Schur functors and the torus grading
(F-016; Macdonald I(4.3), Appendix~A(5.3)--(5.6);
SGA~3 I4.7.3). The two tree copies are auxiliary affine
matrix coordinates, and the coefficient edges encode the
universal linear section of the form bundle. Their determinant
identity is integral in the universal coefficients. Passing
from highest-weight vanishing to the coherent ideal inclusion
uses the characteristic-zero comodule and Schur decomposition.
The inequalities constrain the actual multiplicity test
\(q>K\), at every outer degree, and not only the stronger
sufficient ordinary-Kronecker test. Outside these size bounds,
no evaluation-rank comparison follows from this construction.
